% ============================================================
% Chapter 6: Experimental Design
%   Section: Research Questions and Metrics
% Included from chapters/06-experimental-design/experimental-design.tex
% ============================================================

% Sources: docs/decisions/evaluation.md (the harness, its metrics and
% the stage attribution), docs/experiments.md (discovery, selection,
% admission), docs/roadmap.md "Thesis priorities". Every metric named
% here is computed by a command of the repository; results are in
% Chapter~\ref{ch:results}, not here.

\section{Research Questions and Metrics}
\label{sec:rq-metrics}

Each research question of Section~\ref{sec:research-questions} is
answered by a fixed set of measurements, chosen before the results
were known. Table~\ref{tab:rq-metrics} summarises them; the rest of
this section explains the choices that are not standard.

\begin{table}[ht]
\centering
\small
\begin{tabularx}{\linewidth}{@{}lXXX@{}}
\toprule
& \textbf{What is measured} & \textbf{Metric} & \textbf{Data} \\
\midrule
RQ1 & The work discovery saves, and how well it orders what it
finds & Sources and links read against the every-topic baseline;
precision of the first $k$ candidates shown; AUC of the ranking
score; share of the first $k$ on the topics asked for & Live
candidate rounds, each candidate labelled worth offering or not \\
\addlinespace
RQ2 & Verdict quality, and what retrieval contributes to it &
Accuracy, macro-$F_1$, per-class precision and recall, Cohen's
$\kappa$ against gold, with bootstrap intervals; the same with no
evidence and with the annotator's own evidence; evidence recall and
MRR; the stage each error comes from & The custom set
(Section~\ref{sec:custom-dataset}), natural and synthetic claims
apart; X-Fact as an external benchmark \\
\addlinespace
RQ3 & Agreement between a person and a model making the same
decision & Cohen's $\kappa$ and the confusion matrix; for selection,
precision, recall and Jaccard of the model's picks against the
editor's & Annotator against the AI labelling pass; editor against
the AI selection (blind rounds); annotator against the pipeline \\
\addlinespace
RQ4 & The share of the journalist's work done at a stated quality &
Claim-selector precision; coverage and selective accuracy of the
verdicts; time per article, system against person & All of the
above, joined per article \\
\bottomrule
\end{tabularx}
\caption{Research questions and the measurements that answer them.}
\label{tab:rq-metrics}
\end{table}

\paragraph{RQ1 --- a retriever of news, evaluated as a ranked list.}
The candidates of a discovery round are a ranked list shown to an
editor, so they are evaluated as information retrieval evaluates
one: by the precision of its first $k$ positions and by how well its
score orders relevant above non-relevant items
\cite{jarvelin2002cumulated}, with ``relevant'' meaning \emph{worth
offering}: constructive, on one of the topics asked for, and an
article rather than a page of another kind. The work saved is
counted directly, as the sources and links a scoped round reads
against a round that asks for every topic.

\paragraph{RQ2 --- three conditions, so that retrieval has a number.}
The verifier is run three times over the same claims: with no
evidence at all (the model's own knowledge), with the evidence the
pipeline retrieves, and with the annotator's reference evidence
given directly. The first gap measures what the open-web retrieval
contributes; the second, what better retrieval could still add. This
is the decomposition FEVER and AVeriTeC use when they report
verification with gold and with retrieved evidence
\cite{thorne2018fever,schlichtkrull2023averitec}. Retrieval itself
is measured against the annotator's evidence --- whether its links,
or their domains, are among the candidates, the ranked sources and
the cited ones --- and each wrong verdict is attributed to the stage
of the four-stage architecture \cite{guo2022survey} in which the
trace shows it went wrong. Natural and synthetic claims
(Subsection~\ref{subsec:custom-mutations}) are reported apart, and
evidence published after a claim is flagged rather than counted
\cite{glockner2022missing}.

\paragraph{RQ3 --- agreement, with the annotator as the ceiling.}
Where a person and a model take the same decision on the same input,
their agreement is measured with Cohen's $\kappa$
\cite{cohen1960coefficient}, read with the usual caution about fixed
thresholds \cite{artstein2008intercoder,landis1977measurement}. Three
such decisions exist in the system: which candidates to publish
(editor against the AI selection, with blind rounds so that the
editor cannot be anchored by the model's picks), which label a claim
deserves (the annotator against the AI labelling pass), and the
verdict (the annotator against the pipeline). The annotator's
agreement with their own blind re-label is the ceiling against which
the model's agreement is read; recent work reporting that language
models can match or exceed crowd annotators on some labelling tasks
\cite{gilardi2023chatgpt} is the hypothesis being tested, not an
assumption.

\paragraph{RQ4 --- extent, not replacement.}
The system's answer to RQ4 is a profile, not a verdict: for each step
of the journalist's work, the share the system performs without a
person and the quality at which it does so. The verifier's coverage
(the share of claims it answers with a definitive verdict) is
reported with its selective accuracy (the accuracy on those claims),
because a system that abstains often and is right when it answers
takes over a different share of the work than one that always
answers.
